% ==============================================================================
% PDF-2 — Structural Diagnosability & Phase Logic
% Section 06: Limits and Non-Claims (Reviewer Attack Closure)
%
% Purpose:
%   Close predictable reviewer attacks by making limits explicit: what this paper
%   does NOT do; what cannot be inferred; what would falsify the claims made here.
%
% Target audience:
%   Readers trained in formal methods / systems theory / philosophy of science;
%   hostile reviewers searching for hidden prescriptions or optimization claims.
%
% Language mode:
%   Formal / diagnostic-only. No operational guidance.
%
% Evidence profile:
%   Theorems/definitions internal to the formal setting + ETS.K_EA.v1.1 as an
%   immutable specification boundary. No empirical performance evidence.
%
% Role in argument:
%   Hard boundary layer: prevents misreadings (control, KPI, optimization),
%   clarifies falsifiability scope, and legitimizes STOP as scientific output.
% ==============================================================================

\section{Limits and Non-Claims}
\label{sec:limits-nonclaims}

This document is a \emph{diagnostic} formalization of (i) structural
diagnosability conditions and (ii) phase-logic predicates (including STOP),
\emph{under the assumptions and symbol vocabulary fixed in the formal setting and
by} \texttt{ETS.K\_EA.v1.1}. It does not provide prescriptions, control rules, or
optimization objectives.

\subsection{Non-Claims (Explicit)}
\label{subsec:nonclaims-explicit}

The following statements are \emph{not} asserted anywhere in this paper:

\begin{enumerate}
  \item \textbf{No control or optimization claim.}
  We do not claim the ability to control, optimize, stabilize, or improve an
  enterprise system, nor to select actions that move the system within any phase
  region.

  \item \textbf{No algorithmic procedure claim.}
  We do not provide algorithms, procedures, or operational recipes for computing
  diagnoses or for enforcing constraints in real organizations.

  \item \textbf{No KPI / managerial metric claim.}
  We do not define KPIs, performance targets, maturity models, or managerial
  scorecards. Any numeric illustration elsewhere (if present) is an internal
  demonstration of formal predicates, not a KPI proposal.

  \item \textbf{No causal attribution claim from symptoms to structure.}
  We do not claim that observed incidents uniquely identify internal structure.
  In general, multiple distinct internal structures may be compatible with the
  same observable manifestations.

  \item \textbf{No completeness claim under partial observability.}
  Under incomplete observability, the framework does not guarantee that every
  true internal failure mode is diagnosable from available observations.

  \item \textbf{No universality beyond the formal setting.}
  We do not claim the definitions here are universally adequate for every
  organization without an explicit, justified mapping from that organization to
  the formal objects constrained by \texttt{ETS.K\_EA.v1.1}.
\end{enumerate}

\subsection{Why Control/Optimization Questions Are Ill-Posed Here}
\label{subsec:why-control-illposed}

A recurrent reviewer failure mode is to treat any phase-space representation as
implicitly \emph{actionable}. In this work, that inference is invalid for
structural reasons:

\begin{enumerate}
  \item \textbf{Diagnostic mapping is not a policy.}
  A diagnostic statement has the form ``given assumptions $A$ and observations
  $O$, the set of compatible internal states is $\mathcal{K}_{\text{compat}}$''.
  A control statement requires an additional object: a policy/action space, a
  transition model under interventions, and a goal functional. None of these are
  introduced here, and introducing them would violate the diagnostic-only scope.

  \item \textbf{Underdetermination under partial observability.}
  When observations are incomplete, inverse inference is generally
  non-unique. Without uniqueness, optimization of interventions is not
  well-defined: different compatible internal structures may require incompatible
  interventions.

  \item \textbf{Specification boundary (no new primitives).}
  \texttt{ETS.K\_EA.v1.1} fixes the admissible vocabulary and prevents the covert
  insertion of managerial primitives (objectives, utility, ROI). Therefore,
  ``what should we do?'' is outside the formal language of this paper.

  \item \textbf{Category error: epistemic vs.\ operational.}
  The results here constrain what can be \emph{known} or \emph{ruled out} from a
  given observational interface. They do not specify what can be
  \emph{implemented} in a socio-technical organization.
\end{enumerate}

Thus, control/optimization questions are not ``unanswered''; they are
\emph{ill-posed within the present formalism}.

\subsection{STOP as a Valid Scientific Outcome}
\label{subsec:stop-valid-outcome}

STOP is treated as a first-class outcome because it can be a correct diagnostic
conclusion under the formal semantics:

\begin{itemize}
  \item \textbf{STOP is an epistemic classification, not a recommendation.}
  STOP indicates that, under the formal predicates and available observation,
  the system is in a region where (i) diagnosability fails, (ii) existence
  constraints are violated, or (iii) the formal interface cannot support
  non-degenerate inference.

  \item \textbf{STOP may be the most informative statement possible.}
  In some regimes, ``cannot diagnose within this interface'' is strictly more
  scientifically honest than emitting a spurious positive diagnosis.

  \item \textbf{STOP is falsifiable at the interface level.}
  If an observational regime claimed to be STOP still permits stable, unique
  inference under the same formal assumptions, then the STOP predicate (or its
  prerequisites) is incorrect and must be revised \emph{within the allowed
  vocabulary}.
\end{itemize}

\subsection{Falsifiability Boundaries (What Can Be Refuted)}
\label{subsec:falsifiability-boundaries}

This paper contains two classes of claims, with different falsifiability modes:

\paragraph{(F1) Internal formal claims.}
Statements of the form ``Given Definitions X and Assumptions Y, Proposition Z
holds'' are falsified by a \emph{formal} counterexample: an explicit construction
within the formal setting that violates Z while satisfying X and Y.

\paragraph{(F2) Projection/interface claims.}
Statements relating the formal objects to an empirical enterprise are not
universal; they are conditional on a declared mapping (observational interface).
They are falsified when one can demonstrate at least one of the following:

\begin{enumerate}
  \item \textbf{Spec violation:} the paper implicitly requires primitives not
  permitted by \texttt{ETS.K\_EA.v1.1}, or contradicts an explicit constraint in
  that specification.

  \item \textbf{Incoherent mapping:} the proposed observational interface cannot
  be instantiated without changing the formal meaning of the symbols (i.e., the
  mapping is not semantics-preserving).

  \item \textbf{Interface insufficiency:} for the claimed diagnosability regime,
  the observation channel is provably insufficient (e.g., information loss makes
  the inverse problem non-identifiable) under the assumptions stated.

  \item \textbf{Non-replicability of the predicate evaluation:} if, under the
  same formal assumptions and same declared interface, the predicate evaluation
  is not well-defined (ambiguous semantics), the projection layer is refuted.
\end{enumerate}

\subsection{Hard Limits: What This Framework Cannot Guarantee}
\label{subsec:hard-limits}

Even in its intended scope, the framework cannot guarantee:

\begin{enumerate}
  \item \textbf{Uniqueness of inferred internal state} from observations, unless
  explicit identifiability conditions are met.

  \item \textbf{Completeness of diagnosis} under incomplete observability: some
  internal degradations may remain observationally indistinguishable.

  \item \textbf{Operational feasibility of any intervention,} because
  interventions are not part of the formal language here.

  \item \textbf{Empirical truth of a chosen mapping} from a real organization to
  formal objects: that mapping is an external act, requiring separate evidence
  and validation beyond this paper.
\end{enumerate}

\subsection{What the Reader Should Take Away}
\label{subsec:takeaway}

The correct interpretation is:

\begin{quote}
This paper specifies \emph{conditions under which diagnosis is well-posed} in a
formal sense, and \emph{conditions under which the only honest outcome is STOP}.
It does not prescribe actions, goals, or optimizations.
\end{quote}
